\chapter{Literature Review}
\label{chap:background}

\section{Vision--Language--Action models}
\label{sec:bg-vla}
Classical robot learning trains one policy per task: given enough demonstrations
of ``pick up the red cube'', a policy learns that behaviour and nothing else.
Scaling this way requires a dedicated dataset for every task, object and
environment --- a combinatorial explosion that slow, physical, real-robot data
collection cannot keep up with. Large language models and vision--language models
(VLMs) solved an analogous generalisation problem for text and images by training
one large model on internet-scale data so that it can follow novel instructions.
\emph{Vision--Language--Action} models take that same pretrained vision--language
backbone and extend its output space from text to robot actions, so that broad,
web-scale visual and linguistic understanding can be used to decide what a robot
should do next \parencite{zitkovich2023rt2, sapkota2026vla}.

A VLA is typically built from three parts \parencite{sapkota2026vla}: a
\emph{vision encoder} that turns images into visual features; a \emph{language
model} that reads the instruction and reasons over image and text, usually
inherited wholesale from a pretrained VLM; and an \emph{action decoder} (or action
expert) that turns the fused representation into low-level robot actions. Two
design choices recur across modern VLAs and matter directly for this work:

\paragraph{Action chunking.} Rather than predicting one action per timestep, the
policy predicts a short sequence (a ``chunk'') of future actions in a single
forward pass, then executes some or all of that chunk before re-planning. Action
chunking was introduced by ACT and shown to enable fine-grained behaviours by
modelling temporally consistent action sequences
\parencite{zhao2023act, wang2026vlalimits}. The number of chunk steps actually
executed before re-planning (the execution horizon) has a strong, non-monotonic
effect on success: too short sacrifices smoothness, too long sacrifices
reactivity to a changing environment \parencite{wang2026vlalimits}.

\paragraph{Pretrain-then-post-train.} A VLA's backbone is pretrained once on broad,
non-robot data (or an existing VLM checkpoint), then \emph{post-trained}
(fine-tuned) on a much smaller robot-specific dataset for the target embodiment and
task \parencite{xiang2026posttraining}. This two-stage recipe is the whole point
of using a VLA rather than training from scratch: the model arrives already
knowing what objects and language look like, and need only learn the motor
mapping. This directly motivates the frozen-backbone strategy in
\refChap{chap:methodology}.

\subsection{A short lineage}
\label{sec:bg-lineage}
Following the surveys of \textcite{sapkota2026vla}: early systems
(2022) combined pretrained vision--language embeddings with hand-designed motion
primitives. In 2023, ACT introduced action-chunking transformers
\parencite{zhao2023act}, and RT-2 coined the ``VLA'' framing by fine-tuning a
large pretrained VLM directly on robot trajectories, representing actions as text
tokens so that web-scale knowledge transferred into control
\parencite{zitkovich2023rt2}. In 2024--2025, open-weight and efficiency-focused
VLAs proliferated. $\pi_{0.5}$ demonstrated open-world generalisation to entirely
new homes by co-training on heterogeneous data sources (other robots, subtask
labels, web data), with target-embodiment demonstrations only a small fraction of
training data \parencite{black2025pi05}. SmolVLA, used in this thesis, belongs to
this efficiency-focused generation.

\section{SmolVLA}
\label{sec:bg-smolvla}
SmolVLA is a compact (approximately 450~million parameter) open VLA distributed
through Hugging Face LeRobot, designed for efficient fine-tuning and deployment on
modest hardware \parencite{shukor2025smolvla, cadene2024lerobot}. It combines a
compact pretrained VLM (SmolVLM-2, \cite{marafioti2025smolvlm}) with a flow-matching
action expert. It
processes multi-image and language inputs to produce robot action chunks, and uses
visual-token reduction and a layer-skipping technique that lets the action expert
access features from an intermediate VLM layer; the action expert interleaves
cross-attention and causal self-attention layers
\parencite{kiviranta2025smolvla, shukor2025smolvla}. These properties make it a
capable VLA that fine-tunes comfortably within a 12~GB budget, which is why it is
selected here. The 12~GB ceiling constrained model selection throughout: it ruled out
fine-tuning the multi-billion-parameter families outright, and, as
\refchap{chap:methodology} describes, it also ruled out unfreezing the vision backbone
once three camera streams were in use. No claim is made that SmolVLA is the
\emph{largest} model that would fit; the one larger candidate trialled here,
$\pi_0$, could not be brought up for reasons unrelated to memory
(\refsec{sec:results-journey}), so the question was never settled empirically.

Relative to the discrete-token approach of RT-2 \parencite{zitkovich2023rt2}, which
represents actions as text tokens and decodes them autoregressively, SmolVLA's
flow-matching action head produces continuous action chunks in a single pass; this is
both faster at inference and better suited to the smooth, fine-grained motion that
manipulation requires. Relative to the multi-billion-parameter $\pi_{0}/\pi_{0.5}$ and
GR00T families \parencite{black2025pi05, nvidia2025gr00t}, which target broad
generalisation and (for GR00T) a dual-system reasoning/control split, SmolVLA
sacrifices peak capability for a large reduction in size. That trade is exactly the one
this thesis interrogates: whether the smaller model is nonetheless \emph{sufficient}
for a controlled, language-conditioned manipulation task under a hard hardware ceiling.
The efficiency techniques it employs --- visual-token reduction and intermediate-layer
feature access --- are precisely what allow its compact backbone to be evaluated on
three camera streams within the memory and latency envelope of a consumer GPU.

Kiviranta's bachelor's thesis fine-tunes the same SmolVLA checkpoint on the
low-cost SO-101 arm and reports an 80\% success rate on a single task from around
100 demonstrations, while noting that sequential multi-task fine-tuning caused
catastrophic forgetting and that language conditioning functioned more as a proof
of concept than a fully working instruction-following system
\parencite{kiviranta2025smolvla}. These observations --- both the achievable
success level and the fragility of naive fine-tuning --- are a useful sanity check
for the results in \refChap{chap:results}.

\subsection{A comparison of representative policies}
\label{sec:bg-compare}
\reftab{tab:vlacompare} situates SmolVLA among representative policies relevant to
this thesis. Two axes matter most for the present work: the parameter count (which
determines whether fine-tuning fits a 12~GB GPU) and the action representation
(which affects control smoothness and inference cost). Discrete-token VLAs such as
RT-2 decode actions autoregressively as text, which is expressive but coarse and
slow; flow-matching and diffusion policies such as $\pi_0$/$\pi_{0.5}$ and SmolVLA
emit continuous action chunks, which suit fine manipulation. Among openly available
models, SmolVLA is distinguished less by peak capability than by \emph{efficiency}:
it is the option in this set that fine-tunes and runs within the target budget,
which is precisely why it is selected here.

\begin{table}[ht]
\caption{Representative policies relevant to this thesis. Parameter counts are
approximate and indicative. ACT is included as the non-VLA imitation baseline.}
\begin{center}
\tagpdfsetup{table/header-rows={1}}
\begin{tabular}{lccc}
\toprule
Policy & Approx.\ size & Action representation & Language input \\
\midrule
ACT \cite{zhao2023act}              & $\sim$80\,M   & continuous chunk (from scratch) & no \\
RT-2 \cite{zitkovich2023rt2}        & billions      & discrete text tokens            & yes \\
$\pi_{0.5}$ \cite{black2025pi05}    & $\sim$3\,B    & flow-matching chunk             & yes \\
GR00T N1.6 \cite{nvidia2025gr00t}   & billions      & diffusion (dual-system)         & yes \\
\textbf{SmolVLA} \cite{shukor2025smolvla} & $\sim$450\,M & flow-matching chunk       & yes \\
\bottomrule
\end{tabular}
\end{center}
\label{tab:vlacompare}
\end{table}

\section{ACT: an imitation-learning baseline}
\label{sec:bg-act}
ACT (Action Chunking Transformer) is worth naming precisely because it is used as a
baseline in the preliminary study of \refsec{sec:results-journey}. ACT is a
\emph{from-scratch} imitation-learning policy, \emph{not} a VLA: it has no
pretrained language backbone and cannot accept a free-form text instruction. It is
trained end-to-end on a single task's demonstrations and predicts action chunks
directly from camera images and joint state \parencite{zhao2023act}. It therefore
represents the classical, no-pretraining approach against which a VLA is compared:
ACT shows what a narrow action-chunking policy can achieve on a task, while SmolVLA
shows what the same task looks like with a pretrained vision--language prior.

\section{Imitation learning and behaviour cloning}
\label{sec:bg-il}
Both policies studied in this thesis are trained by \emph{imitation learning}, and
specifically by \emph{behaviour cloning}: a policy $\pi_\theta(\mathbf{a}\mid
\mathbf{o})$ is fitted to reproduce the actions in a fixed dataset of
(observation, action) pairs collected from a demonstrator, here a scripted
inverse-kinematics controller, without any reward signal or environment interaction
during training. The approach dates to ALVINN \parencite{pomerleau1988alvinn}, which
trained a network to steer a vehicle directly from camera input. Behaviour cloning is
attractive for its simplicity and stability and because it requires no reward
engineering, but it has a well-known weakness: \emph{covariate shift}. Because the
policy is only ever trained on states visited by the demonstrator, small errors at run
time move the robot into states that are under-represented in the data, where the
policy is less certain and makes larger errors, which compound --- an effect analysed
formally by \textcite{ross2011dagger}, whose DAgger algorithm addresses it by
iteratively querying the expert on states the learner actually visits. DAgger is not
used here, since the scripted expert could in principle supply such labels but the
training set was fixed in advance. Action chunking mitigates this partially by committing to
temporally consistent multi-step behaviour, which reduces the high-frequency error
accumulation that single-step policies suffer, and is one reason it has become
standard in both ACT and VLAs \cite{zhao2023act, li2025qchunking}.

A complementary remedy is to enrich the demonstration distribution so that it
covers the states a trained policy actually visits, including recoveries from
error. The DAgger family of algorithms formalises this by iteratively rolling out
the current policy, having an expert label the visited states, and retraining;
practical variants collect ``drifted-then-recovered'' demonstrations in which the
imperfect policy acts for a while and a reliable controller then completes the task.
The pilot study underlying this thesis experimented with such recovery data
(\refsec{sec:results-journey}); in the supermarket task, the same intuition appears
as per-episode positional jitter and as the deliberate variety of instruction
paraphrases, both of which broaden the training distribution without an interactive
labelling loop. The relationship between demonstration variety and generalisation is
a recurring theme in the manipulation literature \cite{bharadhwaj2024roboagent,
black2025pi05} and is corroborated first-hand in \refChap{chap:results}.

\section{Robotic grasping and manipulation}
\label{sec:bg-grasp}
Grasping is the core physical skill in this task, and the policy's behaviour is best
understood against the classical grasping literature. Traditional approaches treat
grasp synthesis as an analytic or model-based problem --- computing grasp points that
achieve force closure given the object geometry and friction
\parencite{bicchi2000grasping} --- whereas learning-based approaches predict grasps or
grasping actions directly from sensor data, either by scoring candidate grasps against
learned robustness metrics \parencite{mahler2017dexnet} or, as here, by producing
grasping actions end-to-end as part of a closed-loop policy. The scripted
expert in this thesis is deliberately of the former kind: it uses inverse kinematics
to reach a hand-specified top-down grasp pose, which is reliable precisely because the
object models and poses are known in simulation. The learned policy, by contrast,
must infer an appropriate grasp from images alone, which is why its grasp success
rate (92.5\%) is a meaningful measure of learned competence rather than of scripted
precision.

Two practical grasping considerations shape the design. First, a two-finger parallel
gripper imposes constraints on object geometry: wide or smooth objects are prone to
slipping, which is exactly why the cereal box (widest face) is excluded and why the
water bottle requires a higher, more stable grip point (\refChap{chap:methodology}).
Second, the \emph{approach direction} matters: a top-down grasp followed by a frontal
extraction was adopted because the target shelf bay is enclosed from above, so a
purely vertical lift is infeasible. These constraints are not incidental --- they are
representative of the geometric reasoning any manipulation policy must implicitly
perform, and the per-item variation in success rate in \refChap{chap:results} tracks
them closely.

\section{Language grounding for manipulation}
\label{sec:bg-lang}
For a policy to act on ``pick up the milk carton'', it must \emph{ground} the words
in the visual scene --- associating the token ``milk'' with the corresponding pixels
and, ultimately, with a motor goal. In VLAs this grounding is inherited from the
pretrained vision--language backbone, which has learned joint image--text
representations at scale, rather than being learned from the (small) robot dataset.
This is the mechanism that allows a language-conditioned policy to select the correct
object among distractors and to respond to paraphrases it never saw verbatim. A
central empirical question for any such system, however, is whether the language is
genuinely used or whether the policy has instead latched onto a spurious correlate
(for example, always reaching to a fixed location). \refsec{sec:results-language}
addresses this directly with a controlled test in which only the instruction is
varied, and finds that the reach follows the named item --- evidence of genuine
grounding rather than a memorised trajectory. The fragility of language conditioning
noted by \textcite{kiviranta2025smolvla}, where instruction following functioned more
as a proof of concept, makes this verification particularly important to report
carefully rather than to assume.

\section{Simulation and benchmarks}
\label{sec:bg-sim}
Real-robot data collection is slow, risky and supervision-heavy, so simulation
offers a scalable, safe substitute for benchmarking manipulation policies before
they reach hardware. RoboCasa and RoboCasa365 are large-scale MuJoCo/robosuite
simulation frameworks for training generalist robots on hundreds of household
tasks across thousands of kitchen scenes, arguing that simulation is the practical
way to obtain large, diverse robot datasets \parencite{nasiriany2024robocasa,
nasiriany2026robocasa365}. RoboCasa's reference results use large foundation
policies. The preliminary study in this thesis uses RoboCasa as a stress test of
whether \emph{small}, deployable models can perform its kitchen tasks after
task-specific fine-tuning (\refsec{sec:results-journey}); the negative result there motivates
the controlled supermarket task built in MuJoCo for the remainder of the work.
RoboCasa is demanding for a small model for several compounding reasons: the object
inventory is large and varied, the kitchen scenes are cluttered and visually complex,
the tasks require reasoning about containment (placing \emph{into} a cabinet), and the
reference results are set by a foundation model pretrained on data close to that
distribution. A compact policy fine-tuned from a limited budget must overcome all of
these at once, which is why the sensible response is to reduce the task's difficulty
along each of these axes --- fewer, rigid objects; an uncluttered scene; a simple
open-basket target --- and to establish competence there before reintroducing
complexity. This is the strategy the remainder of the thesis follows. The
same physics-engine family (MuJoCo, via robosuite assets) underpins the custom
environment described in \refChap{chap:methodology} \parencite{todorov2012mujoco}.

The value of simulation is tempered by the \emph{reality gap}: policies trained in
simulation can exploit rendering artefacts, idealised physics or exact object poses
that do not hold on hardware, and so transfer poorly. The dominant mitigation is
\emph{domain randomisation} \parencite{tobin2017domainrand} --- varying textures,
lighting, camera placement, object poses and physical parameters during training so
that the real world appears to the policy as just another variation of the training
distribution. For a VLA there is an
additional, subtler gap: the pretrained vision backbone was trained largely on real
photographs, so synthetic renders are themselves mildly out-of-distribution for it.
Freezing the backbone (\refChap{chap:methodology}) preserves its real-image prior but
cannot adapt it to the synthetic style; a light-weight adaptation such as low-rank
tuning of the vision layers \parencite{hu2022lora} is a plausible middle path,
discussed as future work.
This thesis does not claim sim-to-real transfer and does not employ domain
randomisation, as its aim is to characterise a policy under a controlled,
reproducible simulated task; the pipeline is nonetheless built on the same tooling
(LeRobot, real hardware models) that a subsequent real-robot study would use.

\section{The importance of data and metrics}
\label{sec:bg-data}
A recurring theme in the recent literature is that \emph{data diversity and
evaluation quality} can matter as much as, or more than, model architecture.
$\pi_{0.5}$ attributes its generalisation to data diversity and hierarchical
structure rather than action-model scale \parencite{black2025pi05}; RoboAgent
multiplies a small dataset through semantic augmentation rather than collecting
more data, arguing that data efficiency and diversity outweigh raw dataset size
\parencite{bharadhwaj2024roboagent}. This thesis provides complementary,
first-hand evidence for the same conclusion: as \refChap{chap:results} shows,
fixing the success metric and improving demonstration variety produced larger
gains than model choice at this scale.

\section{Flow-matching action generation}
\label{sec:bg-flow}
The action expert in SmolVLA, like that in $\pi_{0}$/$\pi_{0.5}$, is a
\emph{flow-matching} model rather than a discrete token predictor as in RT-2
\cite{black2025pi05, shukor2025smolvla, zitkovich2023rt2}. Flow matching
\parencite{lipman2023flowmatching} learns a continuous velocity field that transports
a simple noise distribution to the data distribution; at inference the field is
integrated over a small number of steps to sample an output (ten steps in the
configuration used here). For action generation this has two consequences relevant to this
thesis. First, actions are produced as continuous values, which suits smooth
manipulation better than the coarse, autoregressive token decoding of
text-token VLAs. Second, sampling is \emph{stochastic}: repeated queries on the
same observation yield slightly different action chunks, so a policy's success rate
is inherently a distribution rather than a fixed value. This directly motivates the
evaluation protocol adopted in \refChap{chap:methodology}, which reports success
over many trials with confidence intervals rather than from single rollouts.

\section{Evaluation of manipulation policies}
\label{sec:bg-eval}
A recurring methodological hazard in imitation-learning robotics is optimistic or
under-powered evaluation. Because closed-loop rollouts are expensive, results are
often reported over a handful of trials and without confidence intervals, and
success is sometimes measured by a proxy (such as proximity to a goal) rather than
by task completion. Kiviranta reports that a naively defined success metric and
insufficiently varied data inflate SmolVLA's apparent performance
\cite{kiviranta2025smolvla}, and the same effect is documented first-hand in the
pilot study of this thesis (\refsec{sec:results-journey}), where a proximity-based
metric overstated a real 40\% as 80\%. A second hazard is \emph{checkpoint
selection}: training loss and even validation loss continue to fall while
closed-loop task performance may plateau or regress, so selecting a checkpoint by
loss can deploy a worse policy than an earlier snapshot. The evaluation methodology
in this thesis --- adequate sample sizes, Wilson intervals, a task-completion metric
that verifies actual placement, and checkpoint selection by closed-loop success ---
is designed to avoid these hazards, and \refsec{sec:results-checkpoint} demonstrates
the checkpoint effect empirically.

\section{Research gap and positioning}
\label{sec:bg-gap}
The literature establishes that (i) VLAs can transfer web-scale knowledge into
control \cite{zitkovich2023rt2}, (ii) large VLAs generalise to open-world settings
given diverse data \cite{black2025pi05}, and (iii) compact VLAs such as SmolVLA can
be fine-tuned for individual tasks on affordable hardware
\cite{shukor2025smolvla, kiviranta2025smolvla}. What is comparatively
under-reported is a \emph{rigorously evaluated}, end-to-end account of what a
compact VLA can and cannot do under a hard, single-consumer-GPU constraint on a
language-conditioned manipulation task --- including an honest failure analysis and a
disciplined evaluation protocol. This thesis addresses that gap: it fixes the
hardware budget at \SI{12}{GB}, fine-tunes SmolVLA on a controlled supermarket
task, and reports closed-loop success with confidence intervals, a failure
diagnosis, and a language-grounding test. In doing so it also contributes a
first-hand data point --- consistent with $\pi_{0.5}$ and RoboAgent
\cite{black2025pi05, bharadhwaj2024roboagent} --- that at this data and hardware
scale, evaluation quality and data variety dominate model choice.
